% =====================================================================
%  SECTION 11 — Limitations and future work
%  Draft 1. Assembled from caveats stated inside each section, plus
%  those owed to the reader that no individual section carries.
% =====================================================================

\section{Limitations}
\label{sec:limitations}

\subsection{Scope of the construction}

\paragraph{Two dimensions.} The stream-function dual and the Airy
representation~\eqref{eq:airy} are two-dimensional constructions, as is the
Bogner--Fox--Schmit projection used to certify the elastic lower bound. The
three-dimensional case requires a Beltrami-type stress potential together
with a $C^1$-conforming three-dimensional interpolation; the traction
continuity argument of Section~\ref{sec:exact} generalises, but the
implementation does not, and we do not report three-dimensional results.
The same restriction is noted by~\cite{Gaynutdinova2026}, who observe that
the two-dimensional stream-function device has no direct
three-dimensional analogue.

\paragraph{Linear, periodic, small strain.} We treat linear conductivity
and linear plane-strain elasticity on periodic cells with piecewise-constant
coefficients. Nonlinear constitutive behaviour, finite strain, and
non-periodic or statistically defined cells are outside the present
formulation.

\paragraph{Degenerate phases.} The dual side requires a bounded compliance,
so voids ($k \to 0$, $\mathbb{C} \to 0$) and rigid inclusions are not
covered: the complementary energy of the corresponding trial fields is not
finite, and the lower bound degenerates. Porous media must therefore be
treated with a small but nonzero pore phase, and the resulting certificate
bounds that regularised problem rather than the void limit.

\paragraph{Phase-aligned meshes.} Exact integration
(Section~\ref{sec:exact}) requires the certification mesh to resolve the
phase boundaries so that the coefficient is constant per element. For the
pixel-based microstructures used here this is automatic; for microstructures
defined by smooth level sets it requires a conforming mesh, and the
certificate is otherwise exact only for the discretised geometry actually
integrated.

\paragraph{Matrix ordering, not componentwise intervals.} The certificate
is a L\"owner statement,
$C_{\mathrm{low}} \preceq C_{\mathrm{eff}} \preceq C_{\mathrm{up}}$. Scalar
intervals for engineering constants follow from it---for the directional
Young's modulus this is immediate and is what
Section~\ref{sec:applications:design} optimises---but tight componentwise
intervals for individual tensor entries, particularly off-diagonal ones, do
not follow directly from the matrix ordering and require separate
derivation.

\subsection{Where the certificates are too wide to be useful}

Within the controlled generated families of Section~\ref{sec:limits}, the
certificates degrade towards the analytic envelope for corner-singular and
bicontinuous morphologies at high contrast. At contrast $200$ the generated
checkerboard certificate is only $1.2$ times tighter than a bound computable
in closed form from the phase fractions, which is to say it conveys almost
nothing for that case.

Within these controlled numerical cases, the observed widths are consistent
with a trial-field approximation limitation: certified width tracks the
regularity and complexity of the exact corrector. Two targeted remedies
suggest themselves and neither is attempted here. Enriching the trial space with functions carrying the
known corner singularity, in the manner of extended finite element methods,
addresses the singular case directly. A Hashin--Shtrikman polarisation
formulation with a comparison medium is the classical route to high
contrast, although we note it is not obviously the right remedy here, since
a polarisation field would have to represent the same singular behaviour
that the present trial fields represent poorly.

\subsection{Limitations of the amortized operator}

\paragraph{External elasticity scope remains limited.} Section~\ref{sec:amortized:external}
does report full $200/200$ coverage on an independently labelled public
plane-strain dataset after matching the native geometry, phase law, periodic
boundary condition, tensor convention and numerical resolution. That
experiment is a coverage audit, not a competitive benchmark: its median
interval remains $160.9\%$ wide and only $1.05\times$ tighter than the
analytic envelope. We do not report a head-to-head comparison with
label-trained elasticity surrogates on this matched split; that remains
necessary before making any comparative predictive-accuracy claim.

\paragraph{The bottleneck is not compute in the generated family.} Within
the controlled generated family, increasing the training family fivefold
and the optimisation budget ninefold left the held-out median
width unchanged ($28.2\%$ to $28.3\%$), with the monitored gap flat from
step $1200$ onward. We therefore do not claim that additional compute would
tighten these intervals, the ablation of Table~\ref{tab:ablation} identifies network capacity as the limiting factor, since width and depth both help monotonically while
conditioner richness does not. The intervals reported here reflect a
capacity-saturated model rather than an intrinsic bound, and a larger
network would tighten them at proportionate inference cost. Amortized elastic certificates near $28\%$ are useful for
screening but not for final design, where the per-instance route (widths
below $1\%$ in the cases reported) remains preferable.

\paragraph{The geometric-discrimination result is weaker in controlled
elasticity.} The comparison of Section~\ref{sec:tightness:geometry}, in which
microstructures of identical phase fraction receive certificates differing
by more than a factor of two, is compelling for conductivity ($18.3\%$
against $43.2\%$) and much less so for elasticity ($60.1\%$ against
$52.8\%$). We present the conductivity result as primary and report the
elasticity figures alongside; a stronger elastic operator would be required
to make the argument equally forceful in the headline physics.

\paragraph{A matched point-prediction diagnostic is deliberately small.}
On the public DaRUS-5730 elasticity data, we trained three image-based
regressors on the fixed 200-cell training split and evaluated them on the
fixed 100-cell validation split: an unconstrained predictor, an
SPD-by-Cholesky predictor, and a predictor constrained to the analytic
Voigt--Reuss envelope. Their mean relative Frobenius errors were
$23.19\%$, $24.53\%$, and $21.81\%$, respectively. The constrained
predictor was SPD and inside the analytic envelope on all 100 validation
cells; the corresponding rates were $85\%$ and $58\%$ for the
unconstrained predictor, and $100\%$ and $48\%$ for the Cholesky
predictor. These results are a reproducible, matched small-data diagnostic,
not a reproduction of the full-data hyperparameter sweep or reported
accuracy of Keshav, Herb and Fritzen~\cite{Keshav2025}. They also do not
turn a point prediction into a two-sided per-instance certificate. Our
claim remains about the strength and reference coverage of the variational
certificate, not a claim of superior point-prediction accuracy.

\paragraph{The solver audit is conditional.} For the same specified
geometry, phase law, boundary conditions, tensor convention and
discretisation, a reference value falling outside the interval establishes an
inconsistency (Section~\ref{sec:applications:audit}). The diagnostic is
therefore only as sharp as the certificate is tight: it detects an error in
a competing solver only when that error exceeds the certified width. The
case we report was found at contrast $200$, where a finite-volume solution
on a coarse grid fell below the certified lower bound; in the FFT
comparison of Section~\ref{sec:cost} the estimates lay inside the intervals
in every case and no discrepancy was detected. The capability is real but
should not be presented as a routine check.

\subsection{Limitations of the quadrature criterion}

The acceptance thresholds of~\eqref{eq:gate} are empirical, calibrated on
the problems reported here; we do not claim them as universal constants,
and a practitioner adopting the criterion should recalibrate on a case with
a known solution. The bandwidth~\eqref{eq:bandwidth} is defined for
periodic fields, and spectral leakage contaminates the estimate otherwise;
the refinement ratio is free of this restriction and is the diagnostic to
prefer where only one can be afforded. Finally, the criterion is
sufficient for detection but not necessary: we have shown that low headroom
reliably accompanies collocation overfitting in these problems, not that
adequate headroom guarantees a tight certificate.

\subsection{What we have deliberately not claimed}

Two clarifications, since both invite misreading.

First, the certificate is a statement about a model. It asserts that the
effective tensor of the specified periodic, linear, two-phase problem lies
in the reported interval. It is not a prediction about a physical specimen,
and agreement or disagreement with a laboratory measurement would test the
model's idealisations, meaning perfect interfaces, periodicity, plane strain
and linearity, rather than the bound. We have accordingly validated against
closed-form solutions and independent solvers, and, in
Section~\ref{sec:limits:micrograph}, certified the geometry of a real
segmented micrograph without using any measured property. A reader wanting
comparison against measured effective properties will not find it here, and
we regard that as correctly out of scope rather than as an omission.

Second, we do not claim a speed advantage over classical homogenization.
An FFT solve of these cells is roughly fifty times faster per query
(Section~\ref{sec:cost}); the amortized route becomes preferable only in
many-query certification, at approximately $300$ designs, and it is the
guarantee rather than the speed that is the contribution.

\subsection{Future work}

Three directions follow directly. The three-dimensional extension via
Beltrami potentials is the largest and most obviously wanted. Singular
enrichment of the trial space targets the morphologies identified in
Section~\ref{sec:limits} as the principal limitation. Identifying and
removing the bottleneck behind the amortized plateau would make certified
screening tight enough for design rather than triage. Beyond these, the
quadrature criterion of Section~\ref{sec:quadrature} is not specific to
homogenization or to bounded problems, and its systematic characterisation
across collocation-trained physics networks is a separate undertaking.
